\chapter{Computability}
\section{Basic Definitions}

\begin{definition}{Alphabet}{alphabet}
    An alphabet is a set with a finite number of symbols.
\end{definition}

\noindent Usually, an alphabet (\cref{def:alphabet}) is denoted by \(\Sigma\). Now, let us compose words. To do so, we shall use the Cartesian product. As an example, the set of pairs of symbols will be \(\Sigma\times\Sigma=\Sigma^2\). The cardinality of this set is given by the product of the cardinalities of the composing sets: \(|\Sigma\times\Sigma|=|\Sigma|\cdot|\Sigma|\). To represent words with \(n\) symbols, we define the general set of length \(n\) by induction:
\[
    \Sigma^n=\begin{cases}
        \{\varepsilon\} & n=0\\
        \Sigma^{n-1}\times \Sigma & n>0
    \end{cases}
\]
This represents ``The set of all strings on \(\Sigma\) of length \(n\)''.

\noindent To define the set of all words, independent of their length, we take the union of all \(\Sigma^n\) sets:
\[
    \Sigma^0 \cup\Sigma^1 \cup\Sigma^2 \cup\ldots=\bigcup_{n\in\mathbb{N}}\Sigma^n=\Sigma^*
\]

\begin{definition}{Strings on an Alphabet}{strings}
    The strings on an alphabet \(\Sigma\) (\cref{def:alphabet}) are all the strings in its Kleene closure, which is the set:
    \[\Sigma^*=\bigcup_{n\in\mathbb{N}}\Sigma^n\]
\end{definition}

\begin{definition}{Enumeration}{enumeration}
    An enumeration of the strings (\cref{def:strings}) on an alphabet (\cref{def:alphabet}) is a 1:1 and onto function (bijection) that maps the set of natural numbers \(\mathbb{N}\) to the set \(\Sigma^*\).
    \[
        f:\mathbb{N}\rightarrow \Sigma^*
    \]
\end{definition}

\noindent By defining such a bijection between integers and strings, we establish that the set \(\mathbb{N}\) and the set \(\Sigma^*\) have the same cardinality (i.e., they are countably infinite).

\begin{definition}{Decision Function}{decision-function}
    A decision function is a function that takes a string (\cref{def:strings}) in the alphabet \(\Sigma^*\) and maps it to the set \(\{0,1\}\), representing TRUE and FALSE.
    \[
        f:\Sigma^*\rightarrow\{0,1\}
    \]
\end{definition}

\begin{definition}{Partial and Total Functions}{part-total-func}
    A \textbf{total function} maps every element of the domain to exactly one element of the codomain. A \textbf{partial function} is a function that is not total (i.e., it is undefined for at least one element in the domain).
\end{definition}

\begin{definition}{Equal and Distinct Functions}{equal-func}
    Two functions \(f\) and \(g\) are evaluated as equal or distinct based on the following criteria:
    \begin{align}
        f &= g \iff \forall s \in \Sigma^*, \; f(s) = g(s) \\
        f &\neq g \iff \exists s \in \Sigma^*, \; f(s) \neq g(s)
    \end{align}
\end{definition}

\begin{theorem}{Uncountability of Decision Functions}{uncount-dec}
    There is no enumeration (\cref{def:enumeration}) of decision functions (\cref{def:decision-function}).
    \begin{proof}
        Suppose, for the sake of contradiction, that there is a bijection mapping \(\mathbb{N}\) to the set of all decision functions (\cref{def:decision-function}):
        \[
            \mathbb{N}\rightarrow\{f:\Sigma^*\rightarrow \{0,1\}\}
        \]
        Because \(\Sigma^*\) is isomorphic to \(\mathbb{N}\) (as established by \cref{def:enumeration}), this implies a surjective mapping exists directly from any string \(s \in \Sigma^*\) to a decision function \(f_s\):
        \[
            s \mapsto f_s
        \]
        Since this mapping is assumed to be surjective, every possible decision function \(g\) must be accounted for in this enumeration:
        \[
            \forall g \in \{f:\Sigma^*\rightarrow\{0,1\}\}, \quad \exists s \in \Sigma^* \text{ such that } f_s = g
        \]
        Now, let us define a new decision function \(h\) (using Cantor's diagonal argument):
        \[
            h(s) = 1 - f_s(s) \quad \forall s \in \Sigma^*
        \]
        By this definition, it is mathematically true that for every string \(s\):
        \[
            h(s) \neq f_s(s) \implies h \neq f_s
        \]
        Thus, \(h\) is distinct from every \(f_s\) in our enumeration, meaning the function \(h\) is unmapped. This contradicts the assumption that our mapping is surjective. Therefore, a 1:1 mapping (enumeration) of decision functions cannot exist.
     \end{proof}
\end{theorem}
\noindent Thus, there is a fundamental difference in cardinality between the set of strings \(\Sigma^*\) (\cref{def:strings}) and the set of all decision functions \(\{f:\Sigma^*\rightarrow \{0,1\}\}\) (\cref{def:decision-function}). This distinction perfectly illustrates the boundary between \textbf{countable} and \textbf{uncountable} infinite sets.

\begin{definition}{Countable Set}{countable}
    A set \(S\) is \textbf{countable} (or countably infinite) if and only if there exists a bijection (a 1:1 and onto total function) from the set itself to the set of natural numbers \(\mathbb{N}\). Formally, there exists:
    \[
        \exists f: S \rightarrow \mathbb{N} \quad \text{such that } f \text{ is bijective (meaning } f^{-1}: \mathbb{N} \rightarrow S \text{ exists)}.
    \]
\end{definition}

\noindent Based on \cref{def:enumeration} and \cref{def:countable}, we can formally state that the set of strings \(\Sigma^*\) is countable, whereas the set of all decision functions is uncountable.
\begin{definition}{Language}{language}
    A language is a subset, defined by a decision function, of the set of all strings. \(L\subseteq\Sigma^*\) 
\end{definition}
\noindent Furthermore, there is a 1:1 correspondence (bijection) between the set of all decision functions and the power set (the set of all subsets) of \(\Sigma^*\), denoted as \(\mathcal{P}(\Sigma^*)\). 

\begin{definition}{Characteristic Function}{char-func}
    For any subset \(S \subseteq \Sigma^*\) (where \(S \in \mathcal{P}(\Sigma^*)\)), its \textbf{characteristic} (or indicator) \textbf{function} \(\chi_S\) is defined as:
    \[
    \begin{aligned}
        \chi_{S}: \Sigma^* &\rightarrow \{0,1\}\\
        w &\mapsto \begin{cases}
            1 & \text{if } w \in S\\
            0 & \text{if } w \notin S
        \end{cases}
    \end{aligned}
    \]
\end{definition}

\begin{theorem}{Bijection between Decision Functions and the Power Set}{bij-dec-power}
    There is a 1:1 correspondence (bijection) between the set of all decision functions on \(\Sigma^*\) and the power set \(\mathcal{P}(\Sigma^*)\).
    \begin{proof}
        Every decision function \(f\) defines a unique subset of \(\Sigma^*\). We can define a mapping \(C\) from the set of decision functions to the power set \(\mathcal{P}(\Sigma^*)\):
        \[
        \begin{aligned}
            C: \{f : \Sigma^* \rightarrow \{0,1\}\} &\rightarrow \mathcal{P}(\Sigma^*)\\
            f &\mapsto \underbrace{\{s \in \Sigma^* \mid f(s) = 1\}}_{f^{-1}(\{1\})}
        \end{aligned}
        \]
        Conversely, we can define an inverse mapping \(\chi\) using the characteristic function (\cref{def:char-func}) from the power set back to the set of decision functions:
        \[
        \begin{aligned}
            \chi: \mathcal{P}(\Sigma^*) &\rightarrow \{f : \Sigma^* \rightarrow \{0,1\}\}\\
            S &\mapsto \chi_S
        \end{aligned}
        \]
        Because \(C\) and \(\chi\) are inverses of each other (\(C(\chi_S) = S\) and \(\chi_{C(f)} = f\)), they form a bijection. 
    \end{proof}
\end{theorem}

\noindent \Cref{thm:uncount-dec} is a specific application of a much broader mathematical rule, known as Cantor's Theorem:

\begin{theorem}{Cantor's Theorem (General)}{set-pwr-set}
    For any set \(A\), the cardinality of its power set \(|\mathcal{P}(A)|\) is strictly greater than the cardinality of the set itself \(|A|\):
    \[
        |\mathcal{P}(A)| > |A|
    \]
    Furthermore, if \(A\) is a finite set with a cardinality of \(m\) (\(|A| = m\)), then the cardinality of its power set is exactly \(2^m\):
    \[
        |\mathcal{P}(A)| = 2^m
    \]
\end{theorem}

\noindent As an example of defining a set using a finite number of symbols, consider the set of prime numbers \(\mathbb{P}\):
\[
    \mathbb{P} = \{n \in \mathbb{N} \mid n > 1 \land \forall m \in \mathbb{N},\; 1 < m < n \implies m \nmid n\}
\]
An infinite set that can be unambiguously described using a finite string of symbols from a given alphabet is called \textbf{definable}.

\noindent Now, a question arises: \textbf{How many definable sets are there?} \\
They are infinite in number, but strictly \textbf{countable} (since their finite definitions are just strings, and the set of all strings is countable, as established in \cref{def:strings}).

\begin{concept}{Computable Language (Approximate Definition)}{comp-lang-approx}
    We can give an approximate definition of a ``computable language''. A language \(L \subseteq \Sigma^*\) is computable if there exists a program \(P\) (e.g., written in Python, C, etc.) such that for any string \(s\):
    \[
        P(s) = \begin{cases}
            1 & \text{if } s \in L \\
            0 & \text{if } s \notin L
        \end{cases}
    \]
    \textit{Note: The rigorous, formal definition (using Turing Machines) will be introduced later in the course.}
\end{concept}

\begin{definition}{Definable Language}{definable-language}
    A language \(L\) is \textbf{definable} if there exists a first-order predicate \(P\) such that:
    \[
        P(s) \iff s \in L
    \]
\end{definition}

\begin{observation}{Computability vs. Definability}{comp-vs-def}
    \textbf{All computable languages are definable.} The finite source code of the program that computes the language serves as its definition. 
    
    However, \textbf{not all definable languages are computable.} There are languages that can be perfectly defined using formal mathematical logic (such as the Halting Problem) for which no algorithm or program can ever exist.
\end{observation}
\section{A Computational Model: The Turing Machine}

\begin{observation}{Formulas as Strings}{form-strings}
    As a foundational premise: every mathematical formula, no matter how complex, can be unambiguously represented by a single finite string (\cref{def:strings}) of symbols.
\end{observation}

\noindent A program or algorithm can be represented and executed using a theoretical computational model known as a \textbf{Turing Machine}. The machine operates on a one-dimensional tape divided into discrete cells. 

Different definitions of the Turing Machine exist in literature:
\begin{itemize}
    \item \textbf{Tape bounds:} Some definitions require a tape with a defined starting point (e.g., a cell \(0\) exists, making it one-way infinite), while others define the tape as infinite in both directions.
    \item \textbf{Multiple tapes:} Some definitions allow for \(n\)-tapes. However, everything that can be computed on an \(n\)-tape machine can be simulated on a \(1\)-tape machine. Adding more tapes does not increase the computational power (or the class of computable functions) of the machine.
\end{itemize}

\noindent To process numbers and inputs effectively, we designate a \textbf{blank symbol} (usually denoted as \(\sqcup\)) indicating an empty cell, alongside the standard symbols from our alphabet \(\Sigma\) (\cref{def:alphabet}). 

A strict constraint of the machine is that it can only read one cell at a time and has a finite internal memory. To overcome the memory limitation, the machine leverages the infinite tape for storage during computation, which does not reduce its computational capability.

\begin{definition}{Turing Machine States and Transition Function}{tm-trans-func}
    The finite internal memory of the machine is represented by a set of \textbf{states}, denoted as \(Q\). There exists a specific initial state \(q_0 \in Q\) where the computation begins. 

    The computation is driven by a \textbf{transition function} \(F\), which maps the current state and the symbol currently being read to a new state, a symbol to write, and a directional head movement (\(\leftarrow\) for left, \(\rightarrow\) for right):
    \[
    \begin{aligned}
        F: \Sigma \times Q &\longrightarrow \Sigma \times Q \times \{\leftarrow, \rightarrow\}\\
        (\sigma, q) &\longmapsto (\sigma', q', \text{direction})
    \end{aligned}
    \]
    With this simple mechanism, we can represent every computable action, function, or calculation.
\end{definition}

\begin{observation}{The Halt State}{halt-state}
    There is a special state called \textsc{halt} that stops the execution of the transition function. Depending on the formalization, this is handled in the domain/codomain in two mathematically equivalent ways:
    \begin{itemize}
        \item If we consider \(\textsc{halt} \in Q\), the domain excludes it:
        \[
            F: \Sigma \times \left(Q \setminus \{\textsc{halt}\}\right) \longrightarrow \Sigma \times Q \times \{\leftarrow, \rightarrow\}
        \]
        \item If we consider \(\textsc{halt} \notin Q\), it is added to the codomain:
        \[
            F: \Sigma \times Q \longrightarrow \Sigma \times \left(Q \cup \{\textsc{halt}\}\right) \times \{\leftarrow, \rightarrow\}
        \]
    \end{itemize}
\end{observation}

\subsection{Example: Binary Incrementor}
Assume our alphabet (\cref{def:alphabet}) is \(\Sigma = \{\sqcup, 0, 1\}\), where \(\sqcup\) is the default blank symbol. We want to define a transition function (\cref{def:tm-trans-func}) that adds \(1\) to a binary number already written on the tape. 

We define our states as \(Q = \{R, A\}\) (abbreviating \(q_{\text{right}}\) and \(q_{\text{add\_one}}\)). The machine will first move to the far right of the string, and then move left while performing the binary addition (flipping \(1\)s to \(0\)s and carrying the \(1\) until it finds a \(0\) or a blank).

The transition function \(F\) (\cref{def:tm-trans-func}) can be represented using a state table. Note that when transitioning into the \textsc{halt} state (\cref{obs:halt-state}), the computation terminates immediately:

\begin{center}
    \renewcommand{\arraystretch}{1.5}
    \begin{tabular}{r | c c c}
        \textbf{State \textbackslash\ Read} & \(\sqcup\)& \(\mathbf{0}\) & \(\mathbf{1}\) \\
        \hline
        \(R\) \textit{(move\_right)} & \((\sqcup, A, \leftarrow)\) & \((0, R, \rightarrow)\) & \((1, R, \rightarrow)\) \\
        \(A\) \textit{(add\_one)} & \((1, \textsc{halt}, \leftarrow)\) & \((1, \textsc{halt}, \leftarrow)\) & \((0, A, \leftarrow)\) \\
    \end{tabular}
\end{center}

\noindent Nowadays, automata are frequently represented by a directed graph where nodes represent the states (\cref{def:tm-trans-func}) and directed, labeled arcs represent the transitions. The standard notation for an arc label in a Turing Machine is \(\text{Read} / \text{Write}, \text{Direction}\). 

Thus, our binary incrementor can be visually represented by the following state diagram:

\begin{figure}[H]
    \centering
    \begin{tikzpicture}[
        >=Stealth, 
        shorten >=1pt, 
        auto, 
        node distance=4.5cm, 
        semithick,
        every state/.style={thick, fill=gray!10}
    ]
        % Nodes
        \node[state, initial, initial text=start] (R) {\(R\)};
        \node[state] (A) [right=of R] {\(A\)};
        \node[state, accepting] (H) [right=of A] {\textsc{halt}};

        % Edges
        \path[->]
        (R) edge [loop above] node[align=center] {\(0 / 0, \rightarrow\) \\ \(1 / 1, \rightarrow\)} (R)
            edge              node {\(\sqcup / \sqcup, \leftarrow\)} (A)
        (A) edge [loop above] node {\(1 / 0, \leftarrow\)} (A)
            edge              node[align=center] {\(0 / 1, \leftarrow\) \\ \(\sqcup / 1, \leftarrow\)} (H);
    \end{tikzpicture}
    \caption{Incrementor state graph}
    \label{fig:incrementor-graph}
\end{figure}

\subsection{Multiple-Tape Turing Machines}
In the previous definitions, we formalized a Turing machine with a single infinite tape. However, designing algorithms on a single tape can be exceedingly tedious, as the machine must constantly move its head back and forth to keep track of different variables. 

To simplify algorithmic design, we can introduce the \textbf{Multiple-Tape Turing Machine}. A $k$-tape Turing machine has $k$ independent tapes and $k$ independent read/write heads.

\begin{definition}{$k$-Tape Turing Machine Transition Function}{k-tape-tm}
    For a Turing machine with $k$ tapes, the finite set of states $Q$ and the alphabet $\Sigma$ remain conceptually the same, but the transition function $F$ must simultaneously read from and write to all $k$ tapes, as well as move all $k$ heads independently:
    \[
        F: \Sigma^k \times Q \longrightarrow \Sigma^k \times Q \times \{\leftarrow, \rightarrow\}^k
    \]
\end{definition}

While a $k$-tape machine is significantly easier to program, it raises a fundamental question in computability theory: \textit{Does adding more tapes allow the machine to compute functions that a single-tape machine cannot?} The answer is no. 

\begin{theorem}{$k$-Tape Turing Machine Emulation}{k-tape-emulation}
    Any language recognized or function computed by a $k$-tape Turing machine $M_k$ can be recognized or computed by a standard 1-tape Turing machine $M_1$. 
    \begin{proof}
        To emulate $M_k$ on $M_1$, we must represent the contents of all $k$ tapes and the positions of all $k$ heads on a single tape. 
        
        We accomplish this by dividing the single tape of $M_1$ into $2k$ interleaved ``tracks'' (implemented by radically expanding the alphabet). For every tape $i$ in $M_k$, $M_1$ uses two tracks:
        \begin{enumerate}
            \item One track contains the actual symbols of tape $i$.
            \item One track acts as a boolean flag (containing only blanks or a special marker, e.g., $\bullet$) indicating the exact position of head $i$.
        \end{enumerate}
        
        To simulate a single transition of $M_k$, the machine $M_1$ performs the following routine:
        \begin{enumerate}
            \item It scans its entire active tape from the far left to the far right to locate all $k$ head markers.
            \item It stores the $k$ read symbols in its internal finite state memory.
            \item It determines what $M_k$ would do based on its transition function.
            \item It makes a second pass over the tape to update the symbols and move the $k$ head markers according to the transition rule.
        \end{enumerate}
        Because $M_1$ can perfectly replicate every step of $M_k$, they are computationally equivalent. (Note: While they compute the same functions, $M_1$ is quadratically slower, requiring $\mathcal{O}(t^2)$ steps to simulate $t$ steps of $M_k$).
    \end{proof}
\end{theorem}

\subsection{Alphabet Size and Binary Emulation}
Another apparent limitation of our mathematical model might be the size of the alphabet $\Sigma$. Some problems are naturally expressed using massive alphabets (e.g., the entire ASCII or Unicode character sets). Does a machine with a larger alphabet possess more computational power?

\begin{theorem}{Emulation by a Two-Symbol Turing Machine}{two-symbol-emulation}
    Any Turing Machine $M$ operating on an arbitrary finite alphabet $\Sigma$ can be perfectly emulated by a Turing Machine $M_2$ operating on a binary alphabet (e.g., $\Sigma_2 = \{0, 1\}$).
    \begin{proof}
        Let $|\Sigma| = m$ be the number of symbols in the original alphabet. We can encode every unique symbol in $\Sigma$ as a distinct binary string of fixed length $l = \lceil \log_2 m \rceil$. 
        
        The emulating machine $M_2$ will represent each single cell of $M$ using a contiguous block of $l$ cells. To simulate one single step of $M$, $M_2$ must:
        \begin{enumerate}
            \item Read $l$ consecutive binary cells, using its internal states to ``remember'' the sequence of bits.
            \item Once the $l$-bit block is read, $M_2$ knows which original symbol from $\Sigma$ it represents.
            \item $M_2$ transitions to a new state and writes the new $l$-bit encoded symbol back to the tape, sweeping across the $l$ cells.
            \item $M_2$ repositions its head to the beginning of the next adjacent $l$-bit block (either to the left or right, simulating the movement of $M$).
        \end{enumerate}
        Because the mapping between $\Sigma$ and $\Sigma_2^l$ is a bijection, no information is lost, and $M_2$ can perfectly simulate $M$.
    \end{proof}
\end{theorem}

\begin{observation}{Reduction to the Simplest Form}{simplest-tm}
    By combining \cref{thm:k-tape-emulation} and \cref{thm:two-symbol-emulation}, we arrive at a profound conclusion: \textbf{The absolute simplest Turing Machine}—one with only a single tape and a bare-minimum binary alphabet $\{0,1\}$—is mathematically just as powerful as a machine with a million tapes and a billion distinct symbols.
\end{observation}

\section{Universal Turing Machines}
Up to this point, our Turing Machines have been rigidly designed to solve one specific problem (like the binary incrementor in \cref{fig:incrementor-graph}). Their algorithms are ``hardcoded'' directly into their transition functions. If we want to solve a different problem, we must physically design a new machine with a different set of states.

Alan Turing's most revolutionary contribution was realizing that the \textit{description} of a Turing Machine is just a finite set of rules. Since rules can be encoded as strings (\cref{obs:form-strings}), a Turing Machine can read the description of \textit{another} Turing Machine from its tape.

\begin{definition}{\glsfmtfull{utm}}{utm}
    A \textbf{\glsxtrfull{utm}} is a Turing machine $U$ capable of simulating any other Turing machine $M$. 
    
    It takes as input a string $\langle M, w \rangle$, where:
    \begin{itemize}
        \item $\langle M \rangle$ is the encoded string representation (the ``source code'') of the machine $M$'s states and transition function.
        \item $w$ is the input data string that $M$ is supposed to process.
    \end{itemize}
    The machine $U$ will halt and accept/reject the input exactly if and only if $M$ halts and accepts/rejects $w$.
\end{definition}

To build a \gls{utm}, we must formalize an encoding scheme. An excellent practical example of this is a specific \gls{utm} implementation documented in utm.pdf. This \gls{utm} is designed to emulate 3-symbol Turing machines that operate using a blank, a 0, and a 1.

The configuration of the emulated machine $M$ is encoded directly onto the \gls{utm}'s tape in a highly structured format:
\begin{center}
    \texttt{left-tape [ $\hat{M}$ ] current-cell right-tape}
\end{center}
Here, the specific components are defined as follows:
\begin{itemize}
    \item The \texttt{left-tape} and \texttt{right-tape} represent the contents of $M$'s tape to the immediate left and right of the current cell being read.
    \item The $\hat{M}$ inside the brackets encodes the transition rules for $M$, with the current state explicitly marked.
    \item The content of $M$'s current cell is placed immediately following the closing bracket.
\end{itemize}

\paragraph{State and Rule Encoding}
Within the $\hat{M}$ encoding, states are sequentially represented from $q_1$ to $q_m$. To track the computation, state encodings are terminated with a colon (\texttt{:}), except for the currently active state, which is marked with an exclamation point (\texttt{!}). 

Each state's transition rules are broken down into a comma-separated list of three instructions corresponding to reading a blank, a 0, or a 1. Each instruction dictates:
\begin{enumerate}
    \item The target symbol to write.
    \item The direction to shift the head (\texttt{L} or \texttt{R}).
    \item The relative state change, represented by periods (\texttt{.}) for no change, pluses (\texttt{+}) for moving to a higher-numbered state, or minuses (\texttt{-}) for moving to a lower-numbered state.
\end{enumerate}
If a transition rule does not exist for a specific input (indicating a halting configuration), it is simply represented by an empty string. Overall, the spatial complexity (size) of the $\hat{M}$ encoding on the tape is $\mathcal{O}(m^2)$, where $m$ is the total number of states in $M$.

\paragraph{Execution and Performance}
To execute the emulation, the \gls{utm} utilizes an expanded alphabet of sixteen symbols to manipulate the data. The \gls{utm} operates in a continuous ten-step loop: it reads the current cell's symbol, scans left to find the active state marked by the \texttt{!}, locates the appropriate comma-separated rule, marks the rule for execution, writes the new target symbol, shifts the entire $\hat{M}$ encoding left or right to simulate head movement, and updates the state marker before repeating the cycle. The performance of this \gls{utm} is linear relative to the number of steps executed by $M$, and cubic relative to the number of states $m$.

\begin{concept}{The Birth of Software}{birth-of-software}
    The \glsxtrfull{utm} is the mathematical abstraction of the modern stored-program computer (the von Neumann architecture). 
    
    The \gls{utm} acts as the \textbf{hardware} (the \gls{cpu}), which executes a predefined instruction set. The encoded string $\langle M \rangle$ is the \textbf{software} (the program), and $w$ is the \textbf{memory/data}. By proving the \gls{utm} exists, Turing proved that we do not need to build custom hardware for every mathematical problem; we only need to build one flexible machine and feed it different software.
\end{concept}

\subsection{The Church-Turing Thesis}
We should be convinced, by now, that TMs are powerful enough to be a fair computational model, at least on par with any other reasonable definition. We formalize this idea into a sort of ``postulate,'' i.e., an assertion that we will assume to be true for the rest of this course.

\begin{concept}{Church-Turing Thesis}{church-turing-thesis}
    Turing machines are at least as powerful as every physically realizable model of computation.
\end{concept}

\noindent This thesis allows us to extend every result about TMs to every physical computational device.

\subsection{Finding a Non-Computable Function}
At this point, lemma 1 can be given a precise and formal meaning by just replacing the generic word ``algorithm'' with the more precise term ``Turing Machine'': since every TM can be encoded into a string, TMs can only form a countable set, too few to be able to express all decision functions.

Let us now set out to find a function that cannot be computed by any TM (and therefore, by the Church-Turing thesis, by any feasible computational model).

Let us introduce a little more notation. As already defined, the alphabet $\Sigma$ contains a distinguished, ``default'' symbol, which we assume to be $\sqcup$. Before the computation starts, only a finite number of cell tapes have non-blank symbols. Let us define as ``input'' the smallest, contiguous set of tape cells that contains all non-blank symbols at the beginning of the computation.

A Turing machine transforms an input string into an output string (the smallest contiguous set of tape cells that contain all non-blank symbols at the end of the computation), but it might never terminate. In other words, if we see a TM machine as a function from $\Sigma^*$ to $\Sigma^*$ it might not be a total function.

As an alternative, we may introduce a new value, $\infty$, as the ``value'' of a non-terminating computation; given a Turing machine $M$, if its computation on input $s$ does not terminate we will write $M(s) = \infty$.

While TM encodings have a precise syntax, so that not all strings in $\Sigma^*$ are syntactically valid encodings of some TM, we can just accept the convention that any such invalid string encodes the TM that immediately halts (think of $s$ as a program, executed by a UTM that immediately stops if there is a syntax error). This way, all strings can be seen to encode a TM, and most strings just encode the ``identity function'' (a machine that halts immediately leaves its input string unchanged). Let us therefore call $M_s$ the TM whose encoding is string $s$, or the machine that immediately terminates if $s$ is not a valid encoding.

With this convention in mind, we can design a function whose outcome differs from that of any TM. We employ a diagonal technique akin to the proof of Lemma 2: for any string $\alpha \in \Sigma^*$, we define our function to differ from the output of the TM encoded by $\alpha$ on input $\alpha$ itself.

\begin{theorem}{Uncomputable Function $UC$}{uncomputable-uc}
    Given an alphabet $\Sigma$ and an encoding $\alpha \mapsto M_\alpha$ of TMs in that alphabet, the function
    $$
    UC(\alpha) =
    \begin{cases}
        0 & \text{if } M_\alpha(\alpha) = 1 \\
        1 & \text{otherwise}
    \end{cases}
    \quad \forall \alpha \in \Sigma^*
    $$
    is uncomputable.
    \begin{proof}
        Let $M$ be any TM, and let $m \in \Sigma^*$ be its encoding (i.e., $M = M_m$). By definition, $UC(m)$ differs from $M(m)$: the former outputs one if and only if the latter outputs anything else (or does not terminate).
    \end{proof}
\end{theorem}

\noindent What is the problem that prevents us from computing $UC$? While the definition is quite straightforward, being able to emulate the machine $M_\alpha$ on input $\alpha$ is not enough to always decide the value of $UC(\alpha)$. We need to take into account also the fact that the emulation might never terminate. This allows us to prove, as a corollary of the preceding theorem, that there is no procedure that always determines whether a machine will terminate on a given input.

\begin{theorem}{Halting Problem}{halting-problem}
    Given an alphabet $\Sigma$ and an encoding $\alpha \mapsto M_\alpha$ of TMs in that alphabet, the function
    $$
    \text{HALT}(s, t) =
    \begin{cases}
        0 & \text{if } M_s(t) = \infty \\
        1 & \text{otherwise}
    \end{cases}
    \quad \forall(s, t) \in \Sigma^* \times \Sigma^*
    $$
    (i.e., which returns 1 if and only if machine $M_s$ halts on input $t$) is uncomputable.
    \begin{proof}
        Let's proceed by contradiction. Suppose that we have a machine $H$ which computes $\text{HALT}(s, t)$ (i.e., when run on a tape containing a string $s$ encoding a TM and a string $t$, always halts returning 1 if machine $M_s$ would halt when run on input $t$, and returning 0 otherwise). Then we could use $H$ to compute function $UC$.
        
        For convenience, let us compute $UC$ using a machine with two tapes. The first tape is read-only and contains the input string $\alpha \in \Sigma^*$, while the second will be used as a work (and output) tape. To compute $UC$, the machine will perform the following steps:
        \begin{itemize}
            \item Create two copies of the input string $\alpha$ onto the work tape, separated by a blank (we know we can do this because we can actually write the machine);
            \item Execute the machine $H$ (which exists by hypothesis) on the work tape, therefore calculating whether the computation $M_\alpha(\alpha)$ would terminate or not. Two outcomes are possible:
            \begin{itemize}
                \item If the output of $H$ is zero, then we know that the computation of $M_\alpha(\alpha)$ wouldn't terminate, therefore, by definition of function $UC$, we can output 1 and terminate.
                \item If, on the other hand, the output of $H$ is one, then we know for sure that the computation $M_\alpha(\alpha)$ would terminate, and we can emulate it with a UTM $U$ (which we know to exist) and then ``inverting'' the result \`a la $UC$, by executing the following steps:
                \begin{itemize}
                    \item As in the first step, create two copies of the input string $\alpha$ onto the work tape, separated by a blank;
                    \item Execute the UTM $U$ on the work tape, thereby emulating the computation $M_\alpha(\alpha)$;
                    \item At the end, if the output of the emulation was 1, then replace it by a 0; if it was anything other than 1, replace it with 1.
                \end{itemize}
            \end{itemize}
        \end{itemize}
        This machine would be able to compute $UC$ by simply applying its definition, but we know that $UC$ is not computable by a TM; all steps, apart from $H$, are already known and computable. We must conclude that $H$ cannot exist.
    \end{proof}
\end{theorem}

\noindent This proof employs a very common technique of CS, called \textbf{reduction}: in order to prove the impossibility of $\text{HALT}$, we ``reduce'' the computation of $UC$ to that of $\text{HALT}$; since we know that the former is impossible, we must conclude that the latter is too.

\subsubsection*{The Halting Problem for Machines Without an Input}
Consider the special case of machines that do not work on an input string; i.e., the class of TMs that are executed on a completely blank tape. Asking whether a computation without input will eventually halt might seem a simpler question, because we somehow restrict the number of machines that we are considering.

Let us define the following specialized halting function:
$$
\text{HALT}_\varepsilon(s) = \text{HALT}(s, \varepsilon) =
\begin{cases}
    0 & \text{if } M_s(\varepsilon) = \infty \\
    1 & \text{otherwise}
\end{cases}
\quad \forall s \in \Sigma^*.
$$
It turns out that if we were able to compute $\text{HALT}_\varepsilon$ then we could also compute $\text{HALT}$:

\begin{theorem}{Uncomputability of $\text{HALT}_\varepsilon$}{halt-empty}
    $\text{HALT}_\varepsilon$ is not computable.
    \begin{proof}
        By contradiction, suppose that there is a machine $H'$ that computes $\text{HALT}_\varepsilon$. Such machine would be executed on a string $s$ on the tape, and would return 1 if the machine encoded by $s$ would halt when run on an empty tape, 0 otherwise.
        
        Now, suppose that we are asked to compute $\text{HALT}(s, t)$ for a non-empty input string $t$. We can transform the computation $M_s(t)$ on a computation $M_{s'}(\varepsilon)$ on an empty tape where $s'$ contains the whole encoding $s$, but preceded by a number of states that write the string $t$ on the tape. In other words, we transform a computation on a generic input into a computation on an empty tape that writes the desired input before proceeding.
        
        After modifying the string $s$ into $s'$ on tape, we can run $H'$ on it. The answer of $H'$ is precisely $\text{HALT}(s, t)$, which would therefore be computable.
        
        Again, the result was obtained by reducing a known non-computable problem, $\text{HALT}$, to the newly introduced one, $\text{HALT}_\varepsilon$.
    \end{proof}
\end{theorem}

\subsection{Recursive Enumerability of Halting Computations}
Although $\text{HALT}$ is not computable, it is clearly recursively enumerable. In fact, we can just take a UTM and modify it to erase the tape and write ``1'' whenever the emulated machine ends, and we would have a partial function that always accepts (i.e., returns 1) on terminating computations.

It is also possible to output all $(s, t) \in \Sigma^* \times \Sigma^*$ pairs for which $M_s(t)$ halts by employing a diagonal method. Function $\text{HALT}$ is our first example of an R.E. function that is provably not recursive.

Observe that, unlike recursivity, R.E. does not treat the ``yes'' and ``no'' answer in a symmetric way. We can give the following:

\begin{definition}{co-R.E. Decision Function}{co-re-def}
    A decision function $f : \Sigma^* \rightarrow \{0, 1\}$ is co-R.E. if it admits a TM $M$ such that $M(s)$ halts with output 0 if and only if $f(s) = 0$.
\end{definition}

\noindent In other words, co-R.E. functions are those for which it is possible to compute a ``no'' answer, while the computation might not terminate if the answer is ``yes''. Clearly, if $f$ is R.E., then $1 - f$ is co-R.E.

\begin{theorem}{Recursivity vs R.E. and co-R.E.}{re-core-recursive}
    A decision function $f : \Sigma^* \rightarrow \{0, 1\}$ is recursive if and only if it is both R.E. and co-R.E.
    \begin{proof}
        Let us prove the ``only if'' part first. If $f$ is recursive, then there is a TM $M_f$ that computes it. But $M_f$ clearly satisfies both the R.E. definition ($M_f(s)$ halts with output 1 if and only if $f(s) = 1$) and the co-R.E. definition ($M_f(s)$ halts with output 0 if and only if $f(s) = 0$).
        
        About the ``if'' part, if $f$ is R.E., then there is a TM $M_1$ such that $M_1(s)$ halts with output 1 iff $f(s) = 1$; since $f$ is also co-R.E., then there is also a TM $M_0$ such that $M_1(s)$ halts with output 0 iff $f(s) = 0$. Therefore, a machine that alternates one step of the execution of $M_1$ with one step of $M_0$, halting when one of the two machines halts and returning its output, will eventually terminate (because, whatever the value of $f$, at least one of the two machines is going to eventually halt) and precisely decides $f$.
    \end{proof}
\end{theorem}

\noindent Observe that, as already pointed out, any definition given on decision functions with domain $\Sigma^*$ also works on domain $\mathbb{N}$ (and on any other discrete domain), and can be naturally extended on subsets of strings or natural numbers. We can therefore define a set as recursive, recursively enumerable, or co-recursively enumerable.

\subsubsection*{Decision and Acceptance}
In the following, we will use the following terms when speaking of languages.

\begin{definition}{Accept and Decide}{accept-decide}
    If language $S$ is recursively enumerable, i.e. there is a TM $M$ such that $M(s) = 1 \iff s \in S$, then we say that $M$ \textbf{accepts} $S$ (or that it ``recognizes'' it).
    Given a TM $M$, the language recognized by it (i.e., the set of all inputs that are accepted by the machine) is represented by $L(M)$.
    
    If language $S$ is recursive, i.e. there is a TM $M$ that accepts it and always halts, then we say that $M$ \textbf{decides} $S$.
\end{definition}

\noindent In the case of functions transforming strings, we will use the following terms.

\begin{definition}{Compute}{compute}
    If a function $f : \Sigma^* \rightarrow \Sigma^*$ is computable, i.e. there is a TM $M$ that always halts and such that $M(s) = f(s)$, then we say that $M$ \textbf{computes} $f$.
\end{definition}

\noindent We generalize the notion to functions outside the realm of strings by considering suitable representations. E.g., a machine $M$ computes an integer function $f : \mathbb{N} \rightarrow \mathbb{N}$ if it transforms a representation of $n \in \mathbb{N}$ (e.g., its decimal, binary or unary notation) into the corresponding representation of $f(n)$.

Since all representations of integer numbers can be converted to each other by a TM, the choice of a specific one is arbitrary and does not impact on the definition. Therefore, we can resort to unary notation and say that:

\begin{theorem}{Computability of Integer Functions}{comp-int-func}
    A function $f : \mathbb{N} \rightarrow \mathbb{N}$ is computable if and only if there is a TM $M$ on alphabet $\Sigma = \{1, \sqcup\}$ such that
    $$
    \forall n \in \mathbb{N}, \quad M(1^n) = 1^{f(n)}.
    $$
    I.e., the TM $M$ maps a string of $n$ ones into a string of $f(n)$ ones.
\end{theorem}

\subsection{Another Uncomputable Function: The Busy Beaver Game}
Since we might be unable to tell at all whether a specific TM will halt, the question arises of how complex a machine's output can be for a given number of states.

\begin{definition}{The Busy Beaver Game}{busy-beaver}
    Among all TMs on alphabet $\{0, 1\}$ and with $n = |Q|$ states (not counting the halting one) that halt when run on an empty (i.e., all-zero) tape:
    \begin{itemize}
        \item let $\Sigma(n)$ be the largest number of (not necessarily consecutive) ones left by any machine upon halting;
        \item let $S(n)$ be the largest number of steps performed by any such machine before halting.
    \end{itemize}
    Function $\Sigma(n)$ is known as the \textbf{busy beaver function} for $n$ states, and the machine that achieves it is called the Busy Beaver for $n$ states.
\end{definition}

\begin{theorem}{Uncomputability of $S(n)$}{uncomp-sn}
    The function $S(n)$ is not computable.
    \begin{proof}
        Suppose that $S(n)$ is computable. Then, we could create a TM to compute $\text{HALT}_\varepsilon$ (the variant with empty input) on a machine encoded in string $s$ as follows:
        \begin{itemize}
            \item on input $s$, count the number $n$ of states of $M_s$
            \item compute $\ell \leftarrow S(n)$
            \item emulate $M_s$ for at most $\ell$ steps
            \item if the emulation halts before $\ell$ steps, then $M_s$ clearly halts: accept and halt
            \item else $M_s$ takes longer than the BB: reject and halt
        \end{itemize}
        Observe that, by construction, $\Sigma(n) \le S(n)$ (a TM cannot write more than a symbol per step).
    \end{proof}
\end{theorem}

\noindent Given two functions $f, g : \mathbb{N} \rightarrow \mathbb{N}$, we say that $f$ ``eventually outgrows'' $g$, written $f >_E g$, if $f(n) \ge g(n)$ for a sufficiently large value of $n$:
$$
f >_E g \iff \exists N : \forall n > N, \quad f(n) \ge g(n).
$$

\begin{theorem}{$\Sigma(n)$ Outgrows Computable Functions}{sigma-outgrows}
    The function $\Sigma(n)$ eventually outgrows any computable function.
    \begin{proof}
        Let $f : \mathbb{N} \rightarrow \mathbb{N}$ be computable. Let us define the following function:
        $$
        F(n) = \sum_{i=0}^{n} [f(i) + i^2].
        $$
        By definition, $F$ clearly has the following properties:
        $$
        \begin{aligned}
            F(n) &\ge f(n) \quad &\forall n \in \mathbb{N}, \\
            F(n) &\ge n^2 \quad &\forall n \in \mathbb{N}, \\
            F(n + 1) &> F(n) \quad &\forall n \in \mathbb{N}
        \end{aligned}
        $$
        the latter because $F(n + 1)$ is equal to $F(n)$ plus a strictly positive term. Moreover, since $f$ is computable, $F$ is computable too. Suppose that $M_F$ is a TM on alphabet $\{0, 1\}$ that, when positioned on the rightmost symbol of an input string of $x$ ones and executed, outputs a string of $F(x)$ ones (i.e., computes the function $x \mapsto F(x)$ in unary representation) and halts below the rightmost one. Let $C$ be the number of states of $M_F$.
        
        Given an arbitrary integer $x \in \mathbb{N}$, we can define the following machine $M$ running on an initially empty tape (i.e., a tape filled with zeroes):
        \begin{itemize}
            \item Write $x$ ones on the tape and stop at the rightmost one (i.e., the unary representation of $x$);
            \item Execute $M_F$ on the tape (therefore computing $F(x)$ with $C$ states);
            \item Execute $M_F$ again on the tape (therefore computing $F(F(x))$ with $C$ more states).
        \end{itemize}
        The machine $M$ works on alphabet $\{0, 1\}$, starts with an empty tape, ends with $F(F(x))$ ones written on it and has $x + 2C$ states; therefore it is a busy beaver candidate, and the $(x + 2C)$-state busy beaver must perform at least as well:
        $$
        \Sigma(x + 2C) \ge F(F(x)).
        $$
        Now,
        $$
        F(x) \ge x^2 >_E x + 2C;
        $$
        the first inequality comes from our properties, while the second stems from the fact that $x^2$ eventually dominates any linear function of $x$. By applying $F$ to both the left- and right-hand sides, which preserves the inequality sign because $F$ is strictly increasing, we get
        $$
        F(F(x)) >_E F(x + 2C).
        $$
        By concatenating these inequalities, we get
        $$
        \Sigma(x + 2C) \ge F(F(x)) >_E F(x + 2C) \ge f(x + 2C).
        $$
        Finally, by replacing $n = x + 2C$, we obtain
        $$
        \Sigma(n) >_E f(n).
        $$
    \end{proof}
\end{theorem}

\subsection{Reductions}
Note that a few results in the past sections (Theorems 6, 7 and 10) made use of similar arguments: ``If $A$ were computable, then we could use it to solve $B$; however, we know that $B$ is uncomputable, therefore $A$ is too.'' Now we want to formalize such a reasoning scheme.

\begin{definition}{Reduction}{reduction}
    Let $L_1 \subset \Sigma_1^*$ and $L_2 \subset \Sigma_2^*$ be two languages (on possibly different alphabets). A function $f : \Sigma_1^* \rightarrow \Sigma_2^*$ is said to be a \textbf{reduction} from $L_1$ to $L_2$ if
    $$
    \forall s \in \Sigma_1^* \quad s \in L_1 \iff f(s) \in L_2.
    $$
\end{definition}

\noindent Basically, we can use a reduction to transform the question ``Does $s$ belong to $L_1$?'' into the equivalent question ``Does $f(s)$ belong to $L_2$?'' Clearly, to be useful in computability results, $f$ has to be computable.

\begin{definition}{Turing Reduction}{turing-reduction}
    We say that $f : \Sigma_1^* \rightarrow \Sigma_2^*$ is a \textbf{Turing reduction} from $L_1 \subset \Sigma_1^*$ to $L_2 \subset \Sigma_2^*$ if it is a reduction from $L_1$ to $L_2$ and it is computable.
\end{definition}

\noindent If $f$ is a reduction from $L_1$ to $L_2$ we write $L_1 \le_f L_2$. In general, if there is a Turing reduction from $L_1$ to $L_2$ we say that $L_1$ is Turing reducible to $L_2$ and write $L_1 \le_T L_2$.

\begin{theorem}{Reduction Schemes}{reduction-schemes}
    Let languages $L_1$ and $L_2$ and function $f$ be such that $L_1 \le_f L_2$; then
    \begin{enumerate}
        \item if $L_2$ is decidable and $f$ is computable, then $L_1$ is decidable too;
        \item if $L_1$ is undecidable and $f$ is computable, then $L_2$ is undecidable too;
        \item if $L_1$ is undecidable and $L_2$ is decidable, then $f$ is uncomputable.
    \end{enumerate}
    \begin{proof}
        The first point is proven by showing that, if we have a machine for $f$ and a machine for $L_2$, we can build a machine for $L_1$. Let $M_{L_2}$ be a TM that decides $L_2$, and let $M_f$ be a TM that computes $f$. Then the machine $M$ that concatenates an execution of $M_f$ and an execution of $M_{L_2}$, i.e. computes $M(s) = M_{L_2}(M_f(s))$, decides $L_1$ by definition of $f$.
        
        The other two points follow by contradiction.
    \end{proof}
\end{theorem}

\noindent In other words, by writing $L_1 \le_T L_2$ we mean that $L_1$ is ``less uncomputable'' than $L_2$.

\subsubsection*{Consequences of the Halting Problem Incomputability}
If $\text{HALT}$ were computable, we would be able to settle any mathematical question that can be disproved by a counterexample (on a discrete set), such as the Collatz conjecture or Goldbach's conjecture. We would just need to write a machine that systematically searches for one such counterexample and halts as soon as it finds one: by feeding this machine as an input to $H$, we would know whether a counterexample exists at all or not.

More generally, for every proposition $P$ in Mathematical logic we would know whether it is provable or not: just define a machine that, starting from pre-encoded axioms, systematically generates all their consequences (theorems) and halts whenever it generates $P$. Machine $H$ would tell us whether $P$ is ever going to be generated or not.

\section{Rice's Theorem}
Among all questions that we may ask about a Turing machine $M$, some of them have a syntactic nature, i.e., they refer to its actual implementation: ``does $M$ halt within 50 steps?'', ``Does $M$ ever reach state $q$?'', ``Does $M$ ever print symbol $\sigma$ on the tape?''.

Other questions are of a semantic type, i.e., they refer to the language accepted by $M$, with no regards about $M$'s behavior: ``does $M$ only accept even-length strings?'', ``Does $M$ accept any string?'', ``Does $M$ accept at least 100 different strings?''.

\begin{definition}{Property of a Turing Machine}{tm-property}
    A \textbf{property} of a TM is a mapping $P$ from TMs to $\{0, 1\}$, and we say that $M$ has property $P$ when $P(M) = 1$.
\end{definition}

\begin{definition}{Semantic Property}{semantic-property}
    A property is \textbf{semantic} if its value is shared by all TMs recognizing the same language: if $L(M) = L(M')$, then $P(M) = P(M')$.
\end{definition}

\noindent By extension, we can say that a language $S$ has a property $P$ if the machine that recognizes $S$ has it. Finally, we define a property as \textbf{trivial} if all TMs have it, or if no TM has it. A property is \textbf{non-trivial} if there is at least one machine having it, and one not having it.

The two trivial properties are easy to decide, respectively by the machine that always accepts and by the one that always rejects. On the other hand:

\begin{theorem}{Rice's Theorem}{rices-theorem}
    All non-trivial semantic properties of TMs are undecidable.
    \begin{proof}
        As usual, let's work by contradiction via reduction from the Halting Problem.
        
        Suppose that a non-trivial semantic property $P$ is decidable; this means that there is a TM $M_P$ that can be run on the encoding of any TM $M$ and returns 1 if $M$ has property $P$, 0 otherwise.
        
        Let us also assume that the empty language $\emptyset$ does not have the property $P$ (otherwise we can work on the complementary property), and that the Turing machine $N$ has the property $P$ (we can always find $N$ because $P$ is not trivial).
        
        Given the strings $s, t \in \Sigma^*$, we can then check whether $M_s(t)$ halts by building the following auxiliary TM $N'$ that, on input $u$, works as follows:
        \begin{itemize}
            \item move the input $u$ onto an auxiliary tape for later use, and replace it with $t$;
            \item execute $M_s$ on input $t$;
            \item when the simulation halts (which, as we know, might not happen), restore the original input $u$ on the tape by copying it back from the auxiliary tape;
            \item run $N$ on the original input $u$.
        \end{itemize}
        The machine $N'$ we just defined accepts the same language as $N$ if $M_s(t)$ halts, otherwise it runs forever, therefore accepting the empty language. Therefore, running our hypothetical decision procedure $M_P$ on machine $N'$ we obtain ``yes'' if $M_s(t)$ halts (since in this case $L(N) = L(N')$), and ``no'' if $M_s(t)$ doesn't halt (and thus the empty language, which doesn't have the property $P$, is recognized).
        
        Observe that we simply use $N$, which has the property, as a sort of Trojan horse for computation $M_s(t)$.
    \end{proof}
\end{theorem}